\chapter{HASIL DAN PEMBAHASAN}

% TODO: Bab hasil dan pembahasan memuat semua temuan ilmiah yang secara logis
% menerangkan alasan diperolehnya hasil penelitian. Bab diawali dengan
% rancangan eksperimen, dilanjutkan hasil, lalu pembahasan yang menjawab
% rumusan masalah dan tujuan penelitian serta membandingkan hasil dengan
% penelitian terdahulu.

\section{Rancangan Eksperimen}

% TODO: Jelaskan secara detail bagaimana pengujian dilakukan, skenario yang
% digunakan, parameter atau kombinasinya, serta hal-hal yang mengantarkan ke
% hasil eksperimen.

Lorem ipsum dolor sit amet, consectetur adipiscing elit. Sed do eiusmod tempor incididunt ut labore et dolore magna aliqua.

\section{Hasil Eksperimen}

% TODO: Sajikan hasil penelitian secara cermat dalam bentuk tabel, kurva,
% grafik, foto, atau bentuk lain yang lengkap dan jelas (satuan, kondisi
% eksperimen, dan sebagainya). Rujuk setiap tabel/gambar dalam teks dan
% pastikan tabel/gambar tidak muncul sebelum dirujuk.

Lorem ipsum dolor sit amet, consectetur adipiscing elit. Hasil pengujian ditunjukkan pada Tabel~\ref{tab:hasil-eksperimen}.

\begin{table}[H]
\centering
\caption{Hasil Eksperimen}
\begin{tabular}{|l|c|l|}
\hline
\textbf{Aspek} & \textbf{Nilai} & \textbf{Keterangan} \\
\hline
Metrik 1 & 95\% & Lorem ipsum dolor sit amet \\
Metrik 2 & 4.5/5 & Sed do eiusmod tempor \\
Metrik 3 & \checkmark & Ut enim ad minim veniam \\
\hline
\end{tabular}
\label{tab:hasil-eksperimen}
\end{table}

\subsection{Hasil Skenario 1}

% TODO: Uraikan hasil skenario pengujian pertama.

Lorem ipsum dolor sit amet, consectetur adipiscing elit. Sed do eiusmod tempor incididunt ut labore et dolore magna aliqua.

\subsection{Hasil Skenario 2}

% TODO: Uraikan hasil skenario pengujian kedua.

Lorem ipsum dolor sit amet, consectetur adipiscing elit. Sed do eiusmod tempor incididunt ut labore et dolore magna aliqua.

\section{Pembahasan}

% TODO: Bahas hasil penelitian secara sistematis dengan argumentasi rasional,
% baik secara kualitatif maupun kuantitatif. Jelaskan keterkaitan hasil dengan
% rumusan masalah dan tujuan penelitian, serta bandingkan dengan hasil
% penelitian terdahulu (perbaikan, penegasan, atau penolakan interpretasi).

Lorem ipsum dolor sit amet, consectetur adipiscing elit. Sed do eiusmod tempor incididunt ut labore et dolore magna aliqua. Ut enim ad minim veniam, quis nostrud exercitation ullamco laboris nisi ut aliquip ex ea commodo consequat \citep{referensi2}.

\subsection{Perbandingan dengan Penelitian Terdahulu}

% TODO: Bandingkan hasil penelitian Anda dengan penelitian terdahulu.

Lorem ipsum dolor sit amet, consectetur adipiscing elit. Sed do eiusmod tempor incididunt ut labore et dolore magna aliqua \citep{referensi3}.
